\section{Обоснование выбора технологий разработки}

Выбор технологического стека обусловлен требованиями к скорости разработки, безопасности, расширяемости, а также необходимостью интеграции с внешним сервисом расписания. Архитектура приложения включает два независимых сервиса: основной модуль управления задачами и отдельный сервис расписания. Такой подход требует стабильного и хорошо документированного стека технологий.

\subsection{Язык программирования и каркас разработки}

Основной язык разработки --- Python 3.10+. Он обеспечивает быструю разработку и богатую экосистему библиотек для веб-приложений и работы с данными. В качестве веб-фреймворка выбран FastAPI:
\begin{itemize}
    \item высокая производительность благодаря Starlette и асинхронной модели;
    \item строгая типизация и автоматическая валидация данных через Pydantic;
    \item встроенная генерация OpenAPI/Swagger документации для всех эндпоинтов;
    \item естественная поддержка JWT-аутентификации и зависимости через \texttt{Depends}.
\end{itemize}

\subsection{Система управления базами данных и ORM}

В основном сервисе используется реляционная СУБД PostgreSQL. Для доступа к данным применяется ORM SQLAlchemy:
\begin{itemize}
    \item декларативное описание моделей (\texttt{DeclarativeBase});
    \item типизация полей через \texttt{Mapped} и \texttt{mapped\_column};
    \item контроль целостности данных и связей между сущностями;
    \item поддержка транзакций и миграций.
\end{itemize}
Миграции реализованы с помощью Alembic, что позволяет версионировать структуру БД и автоматически применять изменения (\texttt{alembic upgrade head}).

\subsection{Отдельный сервис расписания и MongoDB}

Модуль \texttt{rtu-mirea-schedule} использует MongoDB как документно-ориентированное хранилище. Это оправдано тем, что расписание представляет собой иерархическую структуру (группа → день недели → набор пар), которую удобно сохранять в JSON-подобном формате. Для взаимодействия с MongoDB применяется библиотека \texttt{motor} (асинхронный драйвер).

\subsection{Шаблонизатор и клиентский уровень}

Клиентская часть реализована через серверный рендеринг Jinja2 без SPA-фреймворков. Это снижает сложность проекта и обеспечивает:
\begin{itemize}
    \item быстрый запуск без сложной сборки фронтенда;
    \item возможность интеграции логики с бэкендом в одном проекте;
    \item упрощённую поддержку форм и модальных окон.
\end{itemize}
Для оформления используется собственный CSS (\texttt{app/static/style.css}), что позволяет контролировать визуальные элементы (аватары, карточки дедлайнов, календарь).

\subsection{Инфраструктура и запуск}

Для локальной разработки предусмотрен Docker Compose с PostgreSQL (порт 15432). Это гарантирует единообразную среду выполнения и упрощает настройку БД. Отдельный сервис расписания также запускается в контейнере, что делает архитектуру модульной и независимой.
